\section{Análisis Exploratorio de Datos (EDA)}

En esta etapa se realizó un análisis exploratorio de las variables seleccionadas de tipo \texttt{mean}, con el objetivo de identificar patrones de distribución, relaciones entre variables y diferencias entre las clases \texttt{B} (benigno) y \texttt{M} (maligno).  
Para mantener el análisis enfocado, se excluyeron \texttt{radius\_mean} y \texttt{perimeter\_mean} en algunos apartados, debido a su alta relación con \texttt{area\_mean} y a la redundancia conceptual entre estas variables.

\subsection{Análisis univariado}

Los histogramas muestran que varias variables presentan distribuciones asimétricas y una separación visible entre los diagnósticos benigno y maligno.  
En general, los casos malignos tienden a concentrarse en valores más altos de variables relacionadas con tamaño e irregularidad, especialmente \texttt{area\_mean}, \texttt{concavity\_mean} y \texttt{concave points\_mean}.  
En contraste, los casos benignos se agrupan más hacia valores bajos y con menor dispersión.

Los boxplots refuerzan esta observación.  
Variables como \texttt{concavity\_mean} y \texttt{concave points\_mean} muestran una diferenciación más clara entre clases, lo que sugiere mayor capacidad discriminante.  
Por otro lado, variables como \texttt{texture\_mean}, \texttt{smoothness\_mean}, \texttt{symmetry\_mean} y \texttt{fractal\_dimension\_mean} presentan mayor solapamiento, por lo que su poder de separación individual parece menor.

\subsection{Relaciones bivariadas}

Los scatterplots construidos usando \texttt{area\_mean} como referencia muestran relaciones positivas con varias variables del conjunto.  
En particular, se observa que al aumentar \texttt{area\_mean} también tienden a aumentar variables como \texttt{concavity\_mean}, \texttt{concave points\_mean} y \texttt{compactness\_mean}.  
Esto sugiere que los tumores de mayor tamaño suelen presentar también formas más irregulares.

Además, los casos malignos tienden a concentrarse en regiones de valores más altos, mientras que los benignos permanecen en rangos más bajos.  
Esto indica que las variables morfológicas no solo describen el tamaño del tumor, sino también características asociadas con su grado de irregularidad.

\subsection{Matriz de correlación}

La matriz de correlación evidencia relaciones lineales importantes entre varias variables.  
Se aprecia una fuerte asociación entre las variables relacionadas con tamaño, como \texttt{area\_mean}, \texttt{radius\_mean} y \texttt{perimeter\_mean}, lo cual confirma que contienen información parcialmente redundante.  
También se observa una relación destacada entre \texttt{concavity\_mean} y \texttt{concave points\_mean}, lo que indica que ambas capturan aspectos similares de la morfología del tumor.

En cambio, variables como \texttt{fractal\_dimension\_mean} y \texttt{symmetry\_mean} muestran correlaciones más débiles con el resto, por lo que aportan información complementaria pero menos dominante dentro del conjunto.

\subsection{Análisis multivariado}

El \texttt{pairplot} permite observar simultáneamente la relación entre las variables seleccionadas y el diagnóstico.  
Este gráfico confirma que algunas variables, especialmente \texttt{area\_mean}, \texttt{concavity\_mean} y \texttt{concave points\_mean}, ayudan a separar parcialmente las clases \texttt{B} y \texttt{M}.  
Sin embargo, también se aprecia solapamiento entre observaciones, lo que indica que no existe una separación perfecta usando una sola variable.

En conjunto, el análisis multivariado sugiere que la combinación de varias características morfológicas puede mejorar la capacidad de discriminación entre tumores benignos y malignos.

\subsection{Conclusión del EDA}

El análisis exploratorio muestra que las variables más relevantes para diferenciar entre diagnóstico benigno y maligno son aquellas asociadas al tamaño y a la irregularidad de la masa, en especial \texttt{area\_mean}, \texttt{concavity\_mean} and \texttt{concave points\_mean}.  
También se confirma la existencia de correlaciones fuertes entre algunas variables, lo que sugiere redundancia parcial en el conjunto de datos.

Por esta razón, en etapas posteriores del proyecto se recomienda:
\begin{itemize}
    \item priorizar variables con mayor capacidad de separación entre clases;
    \item revisar la multicolinealidad entre variables altamente correlacionadas;
    \item considerar técnicas de selección de características;
    \item tener en cuenta el desbalance entre clases al construir modelos.
\end{itemize}
